\documentclass[12pt]{article}
\usepackage[margin=1in]{geometry}
\usepackage{booktabs}
\setlength{\parskip}{0.5em}

\title{Reference Tables in Body Text}
\author{legible-pdf synthesized corpus}
\date{}

\begin{document}
\maketitle

\section{First Three Elements}

The lightest three elements of the periodic table appear at the top of the
first row and the start of the second row. Their properties are summarized
in Table~\ref{tab:elements}.

\begin{table}[h]
\centering
\caption{Atomic properties of the first three elements.}
\label{tab:elements}
\begin{tabular}{lccc}
\toprule
\textbf{Element} & \textbf{Symbol} & \textbf{Atomic Number} & \textbf{Atomic Mass (u)} \\
\midrule
Hydrogen & H & 1 & 1.008 \\
Helium   & He & 2 & 4.0026 \\
Lithium  & Li & 3 & 6.94 \\
\bottomrule
\end{tabular}
\end{table}

As shown in Table~\ref{tab:elements}, atomic mass increases monotonically
across these three elements, and the symbol $\mathrm{H}$ corresponds to
the atomic number $Z = 1$.

\section{Astronomical Distances}

Selected nearby objects in the solar system are listed in
Table~\ref{tab:distances}, with distances given in astronomical units.

\begin{table}[h]
\centering
\caption{Approximate distances of selected solar-system objects.}
\label{tab:distances}
\begin{tabular}{lc}
\toprule
\textbf{Object} & \textbf{Distance (AU)} \\
\midrule
Moon    & 0.0026 \\
Mars    & 1.524 \\
Jupiter & 5.203 \\
Saturn  & 9.539 \\
\bottomrule
\end{tabular}
\end{table}

The figures in Table~\ref{tab:distances} reflect mean orbital radii rather
than instantaneous separations, which vary with orbital position.

\section{Glossary}

The document mentions: Hydrogen, Helium, Lithium, H, He, Li,
atomic numbers 1, 2, 3, atomic masses 1.008, 4.0026, 6.94,
Moon, Mars, Jupiter, Saturn, and astronomical units 0.0026, 1.524,
5.203, 9.539. Two tables appear in the document, both with captions.

\end{document}
